% !TeX root = ../../Book2.tex
\section{Hausdorff--Young 不等式}
\label{b2:sec:0402}

变换已经有两个端点估计：可积函数变为有界函数，平方可积函数保持平方范数。中间指数的控制将通过一个连续参数把这两端连起来。这里给出 Fourier 变换所需的完整插值论证，一般算子形式留在\secref{b2:sec:1704}。

\subsection{\texorpdfstring{\(L^1\to L^\infty\)}{L¹→L∞}}
\label{b2:subsec:040201}

由积分三角不等式，\(\|\widehat f\|_\infty\leq\|f\|_1\)。该端点常数为 \(1\)，非负函数在零频率处取等。它强调总质量对每个频率的控制；若只知道 \(L^1\) 范数，就没有进一步统一的频率衰减率。

\subsection{\texorpdfstring{\(L^2\to L^2\)}{L²→L²}}
\label{b2:subsec:040202}

Plancherel 给出另一个端点 \(\|\mathcal F_2f\|_2=\|f\|_2\)。两种定义在交集上一致，才有可能对同一个算子作插值。对有限测度支集的简单函数，它既是绝对收敛积分，也是平方范数变换；下面先在这类共同对象上完成估计。

\subsection{Riesz--Thorin 思想}
\label{b2:subsec:040203}

先准备所需的复参数估计。这里的整函数只用于一个辅助引理；下面出现的函数都是有限指数和，因而可由处处收敛的幂级数直接理解。

\begin{lemma}[三线引理的有界形式]\label{b2:lem:04-three-lines}
设整函数 \(F\) 在闭带 \(0\leq\operatorname{Re}z\leq1\) 上有界，并满足 \(|F(\mathrm{i}t)|\leq1\)、\(|F(1+\mathrm{i}t)|\leq1\)。则 \(|F(\theta)|\leq1\) 对 \(0\leq\theta\leq1\) 成立。
\end{lemma}
\begin{proof}
先回顾幂级数的最大模原理。在某个内点达到模的局部最大值时，把函数乘以常数相位，使中心值为非负实数。幂级数在小圆上的平均等于中心值；而圆上实部不超过该最大值，所以连续的非负差只能处处为零。圆上函数遂为同一常数，逐频率积分又使所有正阶幂级数系数为零。函数在邻域内常值，并由重叠圆盘延拓到连通域。因此非恒定解析函数在紧矩形上的最大模出现在边界。

给定 \(\varepsilon>0\)，考虑 \(F(z)e^{\varepsilon z^2}\)。在两条竖边上，其模至多为 \(e^\varepsilon\)。若 \(|F|\leq M\)，在高度为 \(\pm R\) 的两条横边上，模至多为 \(Me^{\varepsilon(1-R^2)}\)，对充分大的 \(R\) 也不超过 \(e^\varepsilon\)。在矩形上用刚证明的最大模原理，得 \(|F(\theta)|e^{\varepsilon\theta^2}\leq e^\varepsilon\)。最后令 \(\varepsilon\downarrow0\)。
\end{proof}

这一估计是\term[sanxian-yinli]{三线引理}[three-lines lemma]的特殊形式。它把两条边界上的控制传递到中间，下面让边界对应两种不同的积分指数。

\begin{theorem}[Hausdorff--Young]\label{b2:thm:04-hausdorff-young}
设 \(1\leq p\leq2\)，\(q=p'\)，即 \(1/p+1/q=1\)。若 \(f\) 为有界且有限测度支集的简单函数，则
\[
 \|\widehat f\|_q\leq\|f\|_p.
\]
\end{theorem}
\begin{proof}
端点已证，以下设 \(1<p<2\)，并令 \(\theta=2(1-1/p)\)。先归一化 \(\|f\|_p=1\)，再取任意有限测度支集的简单函数 \(g\)，满足 \(\|g\|_p=1\)。记非零复数 \(w\) 的相位为 \(w/|w|\)，零处定义为零。设
\[
 f_z(x)=\operatorname{phase}(f(x))\,|f(x)|^{p(1-z/2)},
 \quad g_z(\xi)=\operatorname{phase}(g(\xi))\,|g(\xi)|^{p(1-z/2)},
\]
零值处仍取零，复幂使用正数 \(|f|\) 的实对数。由于函数各只有有限个非零取值，
\[
 F(z)=\int\!\int e^{-2\pi\mathrm{i}x\cdot\xi}f_z(x)\cdot g_z(\xi)\,\mathrm{d}x\,\mathrm{d}\xi
\]
是有限指数和。其指数中的系数为实数，故在闭带上有界；每个积分的绝对可积性由两个有限测度支集保证。

当 \(\operatorname{Re}z=0\) 时，\(\|f_z\|_1=\|g_z\|_1=1\)，由 \(L^1\to L^\infty\) 端点知 \(|F(z)|\leq1\)。当 \(\operatorname{Re}z=1\) 时，\(\|f_z\|_2=\|g_z\|_2=1\)，由 Plancherel 和 Cauchy--Schwarz 也有 \(|F(z)|\leq1\)。三线引理给出 \(|F(\theta)|\leq1\)；而 \(p(1-\theta/2)=1\)，所以
\[
 \left|\int\widehat f(\xi)g(\xi)\,\mathrm{d}\xi\right|\leq1.
\]
对这些 \(g\) 取上确界即得 \(\|\widehat f\|_q\leq1\)。这里即使尚未知道 \(\widehat f\in L^q\)，对偶范数判别仍合法：先在有限测度区域、有限幅值范围内取其共轭相位乘以 \(|\widehat f|^{q-1}\)，再用简单函数逼近并归一化；得到每个截断的 \(L^q\) 范数不超过 \(1\)，最后用单调收敛。去掉归一化，便得一般结论。
\end{proof}

上述证明体现\term[riesz-thorin-chazhi]{Riesz--Thorin 插值}[Riesz--Thorin interpolation]的方法：变动的是函数幅值的幂，使两个边界恰好落在已知空间。结果称为\term[hausdorff-young-budengshi]{Hausdorff--Young 不等式}[Hausdorff--Young inequality]。本章得到可用常数 \(1\)，不宣称中间指数的最佳常数也等于 \(1\)。

\subsection{\texorpdfstring{\(1\leq p\leq2\)}{1≤p≤2}}
\label{b2:subsec:040204}

\begin{corollary}\label{b2:cor:04-lp-extension}
对每个 \(1<p<2\)，Fourier 变换唯一延拓为连续线性映射 \(\mathcal F_p:L^p\to L^{p'}\)，且 \(\|\mathcal F_pf\|_{p'}\leq\|f\|_p\)。它在 \(L^1\cap L^p\) 上与可积变换一致，在 \(L^2\cap L^p\) 上与 Plancherel 变换一致。
\end{corollary}
\begin{proof}
取简单函数在 \(L^p\) 中逼近，前一定理使变换在 \(L^{p'}\) 中为 Cauchy 列。其极限不依赖逼近列，线性、连续性和唯一性与 \(L^2\) 延拓的证明相同。对交空间中的函数，可先同时截断空间与幅值，再用简单函数同时逼近两种范数。在 \(L^1\) 一端，变换一致收敛；在 \(L^2\) 一端，按平方范数收敛。有限测度区域上的依测度极限唯一，故新延拓与原变换几乎处处一致。
\end{proof}

因此在共同定义域上可以继续使用 \(\widehat f\) 这个记号。若 \(1<p<2\)，变换属于 \(L^{p'}\)，但 \(p'>2\)，不能简单地再对变换使用同一个 Hausdorff--Young 定理来反演；反演的正确表达要看它是否另外满足可积性、平方可积性，或采用后面的分布框架。



\ProseParagraphBreak

一个新定义的必要性，可以从同时不属于两个旧空间的函数看出。固定 \(1<p<2\)，选
\[
 d/2<\alpha<d/p<\beta<d,
\]
并取 \(f(x)=|x|^{-\alpha}\mathbf1_{\{|x|<1\}}+(1+|x|)^{-\beta}\mathbf1_{\{|x|\geq1\}}\)。原点附近由 \(\alpha p<d\) 保证 \(p\) 次可积，无穷远由 \(\beta p>d\) 保证，所以 \(f\in L^p\)。但原点奇性使它不平方可积，无穷远尾部使它不可积。因而它既不能整体代入\chapref{b2:ch:03}的绝对积分，也不能整体代入 Plancherel 定义；Hausdorff--Young 延拓确实增加了对象。

\ProseParagraphBreak

定义这个变换时，只需取在 \(L^p\) 中逼近 \(f\) 的简单函数列。其频率像由不等式组成 \(L^{p'}\) Cauchy 列，不必事先猜测极限的点值公式。若另外将 \(f\) 拆成可积部分与平方可积部分，两个变换的和也应与新定义相容；这种相容性将在下一章的测试配对中统一表达。
